%% 自動生成（phys/gen_phys.py）
\phtocchap{原-11}{質量はエネルギー：相対論と原子核}{太陽は、何をエネルギーのもとにして光っている？}{2,4,5}{3}{5}{460}{C}
\phtocchap{原-12}{こわれる原子核：放射性崩壊と半減期}{大昔の木や骨の年代は、どうやってわかる？}{3,6}{2}{5}{510}{C}
\phtocchap{原-13}{放射線が物にあたる：放射線と物質}{X線写真で、なぜ骨が白くうつる？}{3,5}{2}{5}{552}{C}
\phtocchap{原-14}{放射線をはかる・身を守る：計測と防護}{目に見えない放射線を、どうやってはかる？}{3,6,8}{1}{4}{600}{C}
\phtocchap{原-15}{原子炉の中の中性子：核分裂と原子炉}{原子炉は、どうやって反応を一定に保つ？}{6,8,9}{2}{5}{640}{C}
\phtocchap{原-16}{原子の力とつきあう}{使い終わった核燃料は、どうする？}{}{2}{4}{692}{R}
